%% ============================================================
% f2012_sichuan_sci_q21_locus.tex
%% 2012 四川卷（理科）第 21 题 动点 M 与定点 A(-1,0)、B(2,0)，∠MBA=2∠MAB
%% 来源 PDF：普通高考/2012/2012四川理.pdf 第 2 页（印面 772）第 2 栏
%% 坐标取自原页矢量层（bp），原点为图区左上角 (630.0,515.0)，y 向下。
%% ============================================================
\begin{tikzpicture}[x=1bp,y=-1bp,
  every node/.style={inner sep=0,font=\fontsize{10.5pt}{12.6pt}\selectfont},
  line cap=butt, line join=miter]
  \draw[line width=0.3985bp] (7.656,85.941) -- (135.832,85.941);
  \fill[black] (139.418,85.941) .. controls (138.355,85.742) and (136.629,85.145) .. (135.434,84.449) -- (135.434,87.438) .. controls (136.629,86.738) and (138.355,86.141) .. (139.418,85.941);
  \draw[line width=0.3985bp] (64.328,114.277) -- (64.328,10.188);
  \fill[black] (64.328,6.602) .. controls (64.129,7.664) and (63.531,9.391) .. (62.832,10.586) -- (65.82,10.586) .. controls (65.125,9.391) and (64.527,7.664) .. (64.328,6.602);
  \draw[line width=0.3985bp] (35.992,85.941) -- (106.832,32.105);
  \draw[line width=0.3985bp] (121,85.941) -- (106.832,32.105);
  \node[anchor=base west] at (55.37,9.97) {$y$};
  \node[anchor=base west] at (105.69,31.04) {$M$};
  \node[anchor=base west] at (136.71,96.32) {$x$};
  \node[anchor=base west] at (117.5,98.58) {$B$};
  \node[anchor=base west] at (29.3,98.87) {$A$};
  \node[anchor=base west] at (51.74,100.41) {$O$};
\end{tikzpicture}
